%====================
% INRIA
%====================
\subsection{{AI Researcher \hfill Oct.\ 2024 --- Oct.\ 2027}}
\subtext{Institut national de recherche en sciences et technologies du numérique (INRIA) \hfill Villeneuve-d'Ascq, France}
\begin{zitemize}
    \item Improving code completion and code generation using LLMs, specifically targeting the Pharo programming language, which has limited training data.
    \item Developing techniques for code completion, type inference, and deployment in Pharo’s IDE, with a focus on runtime performance.
\end{zitemize}

%====================
% INERIS
%====================
\subsection{{AI Engineer \hfill Sept.\ 2023 --- Sept.\ 2024}}
\subtext{Institut national de l'environnement industriel et des risques (INERIS) \hfill Verneuil-en-Halatte, France}
\begin{zitemize}
    \item Developed \textit{INERIS-IA}, a tool to classify textual documents according to INERIS’s strategic goals using ML and NLP techniques.
    \item Created Boolean queries for document retrieval and improved corpus quality through document similarity and keyword extraction.
\end{zitemize}

%====================
% LIPN
%====================
\subsection{{Intern --- AI Researcher \hfill June\ 2023 --- Aug.\ 2023}}
\subtext{Laboratoire d'Informatique de Paris Nord (LIPN) \hfill Villetaneuse, France}
\begin{zitemize}
    \item Compared sampling techniques for probabilistic planning, particularly for generating literary narratives.
    \item Evaluated methods including the Score Function Estimator (SFE) and advanced techniques such as Gumbel-Softmax, assessing their effectiveness in producing coherent and creative stories.
\end{zitemize}